\usepackage{needspace}

\newenvironment{AutoBibliography}
  {\begin{small}}
  {\end{small}}

% Minimum number of lines that must remain before starting an entry.
\newcount\AutobibNeedlines
\AutobibNeedlines=5

% Hanging indentation of ordinary bibliography entries.
\newlength{\ABindent}
\setlength{\ABindent}{0.05\linewidth}

% Additional indentation of subsequent paragraph first lines.
\newlength{\AutobibParindent}
\setlength{\AutobibParindent}{1.5em}

% Width reserved for the number in a numbered bibliography.
\newlength{\AutobibNumberWidth}
\setlength{\AutobibNumberWidth}{5ex}

\newcommand{\Autobibentry}[1]{%
  \par
  \Needspace{\AutobibNeedlines\baselineskip}%
  \begingroup
    \emergencystretch=2em
    \tolerance=1000
    \leftskip=\ABindent
    \parindent=\AutobibParindent
    \noindent
    \hspace*{-\ABindent}%
    #1%
    \par
  \endgroup
  \addvspace{1.0ex}%
}

\newcommand{\Autobibtarget}[1]{\phantomsection#1}
\newcommand{\Autocolbibnumber}[1]{%
  \par
  \Needspace{\AutobibNeedlines\baselineskip}%
  \begingroup
    \emergencystretch=2em
    \tolerance=1000
    \leftskip=\AutobibNumberWidth
    \parindent=\AutobibParindent
    \noindent
    \hspace*{-\AutobibNumberWidth}%
    \makebox[\AutobibNumberWidth][r]{#1~~}%
    \ignorespaces
}
\newcommand{\Autocolbibentry}[1]{%
    #1%
    \par
  \endgroup
  \addvspace{1.0ex}%
}
\newcommand{\Autobibref}[1]{#1}
\providecommand{\AutobibLink}[1]{#1}

\newcommand{\pseudodoi}[1]{#1}
